\clearpage
\fancyhead[L]{\small STEINHEIL. PRISMENPHOTOMETER.}
\SourceLeaf{654}
\begin{center}\small Göttingische gelehrte Anzeigen. Stück 34. 35. Seite 329..337. 1835 März 5.\\[3pt]\rule{.72\linewidth}{.3pt}\end{center}

Göttingen. Zur Beantwortung der auf den November 1834 von der Mathematischen Classe der Königlichen Societät aufgegebene Hauptpreisfrage, deren Termin aber nach St. 149 dieser Anz. vom v. J. bis Ende December verlängert war, waren drei Concurrenzschriften eingelaufen, eine in lateinischer Sprache mit dem Motto: Opinionum commenta delet dies, naturae iudicia confirmat; die zweite in deutscher Sprache mit der Aufschrift: Suum cuique; die dritte gleichfalls deutsch mit den Worten: Nur gleichartige Eindrücke sind vergleichbar.

Die Abhandlung Nr. 2, mit der Aufschrift: Suum cuique, enthält nur die Meinungen ihres Verf. über die Bildung und Naturbeschaffenheit der Himmelskörper, und gar nichts, was auf die Lösung der von der Societät gestellten Aufgabe Bezug hätte. Eine besondere Beurtheilung jener Meinungen ist daher unnöthig, da solche mit der Preisfrage in gar keinem Zusammenhange stehen.

Der Verf. der Schrift Nr. 1, Opinionum commenta u. s. w. hat hingegen die Frage richtig aufgefasst, einen Apparat zur Vergleichung der Lichtstärke zweier Sterne angegeben und ausführen lassen, auch einige Versuche der Anwendung auf wirkliche Lichtmessungen mitgetheilt. Das Instrument ist ein Fernrohr mit solchen Vorrichtungen, dass beide Sterne zugleich im Felde neben einander gesehen werden können, der eine direct, der andere durch Reflexion. Letztere wird durch einen vor dem Objectiv angebrachten Spiegel bewirkt, der sich in die dem Winkelabstande beider Sterne entsprechende Neigung gegen die Gesichtslinie durch Drehung um eine die Gesichtslinie rechtwinklicht schneidende Axe bringen lässt; der äussere Rand des Spiegels fällt mit dieser Drehungsaxe zusammen, daher der Spiegel in jeder Lage die Hälfte des Spiegels für directes Licht verstattet. Es ist nun aber noch unmittelbar vor dem Objectiv eine halbkreisförmige Blendung angebracht, welche nur die Hälfte des Objectivs offen lässt und ganz herumgedreht werden kann. Die Grösse dieser Drehung wird auf einem eingetheilten Ringe (so wie die Grösse der Spiegeldrehung auf einem Gradbogen) gemessen. Steht der Index des Ringes auf dem Nullpunkt, so kommt gar kein directes, nach einer halben Umdrehung hingegen kommt gar kein reflectirtes Licht in das Fernrohr: bei jeder Zwischenlage theilt sich das reflectirte und das directe Licht im Verhältniss der Abweichung von jenen beiden Stellungen in die offene Hälfte des Objectivs. Man übersieht so leicht, dass, wenn man durch Drehung der Objectivblendung bewirkt hat, dass beide Sterne gleich hell erscheinen, sich, vorbehältlich eines noch unbekannten von der Schwächung des Lichts durch die Reflexion abhängigen Factors, das Verhältniss der Lichtstärke beider Sterne berechnen lässt: dieser unbekannte Factor wird gefunden oder eliminirt durch Zuziehung einer zweiten Beobachtung, wobei blos die Sterne vertauscht werden. Für gewisse Fälle hat der Verf. noch einen zweiten Spiegel beigefügt, so dass der eine Stern durch doppelte Reflexion gesehen wird, was übrigens in der Methode keinen Unterschied macht. Die Bequemlichkeit des Gebrauchs wird durch ein parallactisches Stativ sehr erhöht.

Man muss bedauern, dass der späte Empfang dieses Instruments aus den Händen des Verfertigers den Verfasser gehindert hat, eine durchgreifende Prüfung durch zahlreiche Messungen auszuführen. Er hat das Lichtverhältniss von sieben Sternpaaren, zusammen aus nur 44 Beobachtungen, die jedoch nur summarisch angezeigt werden, bestimmt. Die Resultate, die zuerst gesetzten Sterne jedesmal als Einheit betrachtet, sind folgende:
\SourceLeaf{655}
\begin{center}
\begin{tabular}{lc}
Sterne & Lichtverhältniss\\[3pt]
Rigel, Procyon & $0.8501$\\
Rigel, $\beta$ kl. Hund & $0.1258$\\
Sirius. Rigel & $0.2875$\\
Sirius, Procyon & $0.2756$\\
Procyon, Regulus & $0.3781$\\
Procyon, Nordstern & $0.4369$\\
Regulus, Nordstern & $0.5720$
\end{tabular}
\end{center}

Die Höhen der Sterne, oder die Grössen, wovon sie abhangen, fehlen. Die wahrscheinlichen Fehler dieser Bestimmungen, so weit sie aus der Vergleichung der einzelnen Beobachtungen unter sich festgesetzt werden können, würden nach den Anführungen des Verf. zwischen $\frac{1}{16}$ und $\frac{1}{64}$ des Ganzen schwanken. Vergleicht man nun aber die erste, dritte und vierte Bestimmung unter sich, so zeigt sich die Nothwendigkeit viel stärkerer Correctionen, und die drei letzten Bestimmungen lassen sich gar nicht vereinigen. Der Verf. gesteht selbst, dass er diesen Widerspruch nicht zu erklären wisse, und wenn man gleich hoffen muss, dass es ihm in Zukunft nach viel umfassenderen Versuchen gelingen werde, die Quelle solcher Fehler aufzufinden, so bleibt doch gegenwärtig die Tauglichkeit des Apparats zur Messung der Helligkeit leuchtender Punkte noch unverbürgt.

Der Verf. der dritten Abhandlung mit dem Motto: \emph{Nur gleichartige Eindrücke sind vergleichbar}, hat zwei ganz verschiedene Apparate angegeben und ausgeführt: den einen nennt er den Ocularapparat, den andern das Prismenphotometer. Obwohl beide zu dem vorgegebenen Zweck angewandt werden können, so ist doch eigentlich der erstere weniger zur Vergleichung der Lichtstärke leuchtender Punkte, als zur Vergleichung der specifischen Helligkeit ausgedehnterer Flächen, z. B. des Himmelsgrundes, bestimmt, und es wird daher hinreichen, hier nur die Hauptmomente des zweiten Apparats anzugeben. Der Grundgedanke für dieses Instrument ist die bekannte Erfahrung, dass ein Stern, welcher dem unbewaffneten Auge, oder in einem zum deutlichen Sehen gestellten Fernrohr wie ein untheilbarer leuchtender Punkt erscheint, sich in ein kreisförmigs Bild ausbreitet, wenn man dem Ocular eine andere Stellung gibt, als das deutliche Sehen erfordert. Dieses Bild ist desto grösser, aber eben deshalb in seinen Theilen desto lichtschwächer, je weiter das Ocular von seiner Normalstellung absteht. Für ungleich helle Sterne muss man daher das Ocular in ungleiche Entfernung von der Normalstellung bringen, um die Bilder in gleicher Flächenhelligkeit erscheinen zu lassen. Es lässt sich so die Lichtstärke zweier Sterne schon einigermassen vergleichen, wenn man undeutliche Bilder von ihnen \emph{nach einander} beobachtet, ihre Flächenhelligkeit, so viel der Gedächtnisseindruck verstattet, gleich macht und die entsprechenden Ocularstellungen abmisst. Natürlich erwartet man von einem so rohen Verfahren wenig Genauigkeit und findet sich daher überrascht, dass die von dem Verf. angeführten Versuche eine doch viel grössere Uebereinstimmung darbieten, als man hätte erwarten mögen: dies erweckt schon ein günstiges Vorurtheil für den von dem Verf. kunstreich angeordneten Apparat, womit man derartige Bilder zweier Sterne \emph{zugleich} sehen und zu gleicher Flächenhelligkeit bringen kann.

Das Objectiv ist in zwei gleiche Hälften zerschnitten, die sich nicht neben einander, wie am Heliometer, sondern längs ihrer gemeinschaftlichen Axe, jede für sich, verschieben lassen. Die Mitte der Verschiebungen, die durch Scalen an der Aussenseite des Rohrs scharf gemessen werden, entspricht, wenn die Ocularröhre ganz eingeschoben ist, ungefähr derjenigen Stellung gegen letzteres, die zum deutlichen Sehen erfordert wird. Die beiden Objectivhälften erhalten ihr Licht durch Spiegel, deren reflectirende Flächen $45^\circ$ gegen die Axe des Rohrs geneigt sind, und von denen der eine (vom Objectiv weiter abstehende) um diese Axe messbar gedreht werden kann. Diese Axe ist also beim Beobachten zweier Sterne immer gegen den einen Pol des sie verbindenden grössten Kreises zu richten. Die \SourceLeaf{656}Spiegel selbst sind Glasprismen, in welche das Licht senkrecht einfällt, und senkrecht aus ihnen
austritt. Zwischen den Objectivhälften und den zu ihnen gehörenden Prismenspiegeln sind
Diaphragmen angebracht, die durch zwei Schieberpaare gebildet werden: jedes Schieberpaar wird
durch Eine Schraube mit entgegengesetzt geschnittenen Gewinden so bewegt, dass die Mitte der
Hypotenuse des zu einem grössern oder kleinern rechtwinkligen Dreiecke sich bildenden Diaphragma
unverrückt bleibt.

Vermöge dieser Einrichtung sieht man bei gehöriger Stellung des Rohrs und der Spiegel zwei Sterne
zugleich, und zwar jeden wie eine rechtwinklige Dreiecksfläche, wenn die Objectivhälften von der
Normallage zum Ocular abweichen: von dieser Abweichung hängt sowohl die scheinbare Grösse des
Dreiecks, als dessen Flächenhelligkeit ab, aber jene zugleich mit von der Diaphragmenöffnung, diese
von der eigenthümlichen Helligkeit jedes Sterns: man kann daher durch Aenderung der einen
Abweichung die Flächenhelligkeiten beider Bilder, und wenn man will, durch Abänderung einer
Diaphragmenöffnung, auch ihre Grösse, zur Gleichheit bringen. Dass so das Verhältniss der
Lichtstärke zweier Sterne gefunden, und dabei auch etwaige Ungleichheiten in den Objectivhälften
und Prismenspiegeln durch umgekehrte Combination eliminirt werden können, bedarf nun keiner
weitern Ausführung.

Der Verf. hat seinen Apparat einer strengen Prüfung unterzogen, aber geflissentlich nicht an
Sternen, sondern an künstlich hervorgebrachten sternähnlich leuchtenden Punkten. Diese künstlichen
Sterne erhielt er durch den Reflex des Tageslichts von zwei nahe gleichen gut polirten Stahlkugeln,
etwa $\frac{1}{2}$ Zoll im Durchmesser. Das Tageslicht, für beide Kugeln von einerlei Stelle des
Himmelsgrundes herrührend, gelangte zu den Kugeln durch kreisrunde Blendungen von verschiedener
Weite, und es war Sorge getragen, dass kein fremdes Licht weder die Kugeln noch das Auge des
Beobachters treffen konnte. Es wurden überhaupt vier Blendungen gebraucht, die engste 7, die
weiteste 20 Linien im Durchmesser; durch die sechs verschiedenen Combinationen konnte man also
künstliche Sterne von sechs verschiedenen Lichtverhältnissen erhalten; die grösste Ungleichheit,
wie 1 zu 8, entspricht nach des Verf. eigenen Untersuchungen nahe dem Mittelverhältnisse zweier
Sterne, die um zwei Ordnungen von einander abstehen. Diese künstlichen Sterne erscheinen wirklich
ganz ähnlich, aber ohne den Wechsel und das Wallen, wodurch die Beobachtungen wirklicher Sterne
oft so unsicher werden: überdies hatten sie den höchst wichtigen Vorzug, dass ihr
Helligkeitsverhältniss aus den Blendungsöffnungen a priori bekannt war. Der Verf. theilt die
grosse Zahl der Messungen ihrer Lichtstärke mit dem Prismenphotometer im ausführlichen Detail mit,
ohne diejenigen zu verschweigen, bei welchen sich anfangs einige Unregelmässigkeiten zeigten, deren
Ursachen jedoch entdeckt und weggeräumt wurden. Der wahrscheinliche Fehler Einer Vergleichung
ergibt sich aus der Gesammtheit der Messungen als $\frac{1}{32}$
der ganzen Helligkeit, diese möge gross oder klein sein, und die Verhältnisse der verschiedenen
künstlichen Sterne zeigen eine vollkommen befriedigende Uebereinstimmung mit den
Blendungsöffnungen.

Die Tauglichkeit des Apparats zu scharfer Vergleichung der Helligkeit leuchtender Punkte ist
hierdurch auf eine genügende Art erwiesen, und wenn man auch ungern Anwendungen auf wirkliche
Sterne vermisst, so hat man doch Grund genug, auch bei diesen befriedigende Resultate zu erwarten,
wenn man nur, wie der Verf. mit Recht verlangt, die Beobachtungen auf besonders günstige
atmosphärische Zustände beschränkt, wo man, bei der leichten Handhabung des Instruments, in
wenigen Stunden mehr ausrichten wird, als unter ungünstigen Umständen an vielen Tagen. Uebrigens
enthält die Abhandlung noch manche andere photometrische Untersuchungen und Ansichten von
bedeutendem Interesse, die jedoch, als zur Hauptsache nicht wesentlich nothwendig, hier mit
Stillschweigen übergangen werden können. Einige Anwendungen der Wahrscheinlichkeitsrechnung im
letzten Abschnitt würden einer Berichtigung bedürfen, was jedoch für den Hauptgegenstand selbst
ganz unwesentlich ist.

Endlich kann noch bemerkt werden, dass das Prismenphotometer, obwohl auf ein ganz anderes
\SourceLeaf{657}Princip gegründet, als das der Abhandlung Nr. 1 zum Grunde liegende, doch zugleich die
Möglichkeit darbietet, Sterne nach dem andern Princip zu vergleichen, nemlich durch zugleich
erscheinende \emph{deutliche} Bilder bei messbar verengter Objectivöffnung, und dass selbst bei
dieser Beobachtungsart, welche übrigens der Verf. nach seinen Erfahrungen für verwerflich hält,
die Einrichtung des Prismenphotometers Vorzüge vor der bei Abhandlung 1. beschriebenen haben
würde.

Da die Abhandlung 3. die Aufgabe am vollkommensten und auf eine solche Art gelöst hat, dass ein
schätzbarer Fortschritt in diesem Theile der praktischen Astronomie dadurch begründet wird, so hat
die königl. Societät ihr den Preis, der Abhandlung 1. hingegen, die ebenfalls sehr verdienstvoll
ist, das Accessit zuerkannt.

Der Verfasser der gekrönten Abhandlung ist, nach dem in der öffentlichen Sitzung der Societät vom
14. Februar entsiegelten Zettel,

\begin{center}
Dr. \textsc{Steinheil} in München.
\end{center}

Der Zettel zu der Abhandlung Nr. 2. wurde in derselben Sitzung uneröffnet verbrannt.
